% Transcription — fonds Alexandre Grothendieck, shelfmark 34, batch 8 (pages 141-160).
% Established from the facsimile archives/batches/34/batch-08.pdf (a mirror of
% https://grothendieck.umontpellier.fr/34.pdf), per the skill
% transcribe-grothendieck.
% Manuscript draft of SGA 7, his pagination V.14-V.33 (one exposé, numbered V):
% end of §3 (global vanishing of vanishing cycles), §4 the local situation and
% the five schemes V_s, V, V_eta, V_etabar, Vbar with their exact triangles,
% §5 application of the local depth theorems (5.2-5.7, local acyclicity and
% 1-asphericity, conjectures), §6 an isolated non-smooth point (6.1-6.5:
% the triangle of u, beta, local cohomology with supports in V_s).
% All twenty pages are in his hand and transcribed.
% Pass: Opus 5.5 (claude-opus-5-5), 2026-10-03 — first pass, unchecked against the pages by a human.
%
\documentclass[11pt,a4paper]{article}
\input{../preamble/grothendieck}

\folder{34}
\batch{8}
\pages{141}{160}
\foldertitle{SGA 7 : notes manuscrites (s.d.).}
\dating{[vers 1967-1973]}
\watermark{Édition de démonstration}

\begin{document}

\page{141}
\note{Le feuillet porte la pagination de l'auteur « V.14 » ; la remarque 3.2 prolonge le §3 commencé avant ce lot.}

\emph{Remarque} 3.2. a) Lorsque $X_s$ est \add{localement} une intersection complète,
\struck{alors} et purement de dimension $n$, alors \struck{l'hyp. p\ill{}}
\struck{les hypothèses} il en est de même de $X_{\bar\eta}$, et les hypothèses (3.1.1) et
(3.1.2) se vérifient, avec $m = n$. \struck{Alors} Nous verrons cependant plus bas
(au \struck{\ill{}} §~5) que l'on peut (sous réserve de conjectures raisonnables
\struck{dans} (3.3)) \uncertain{préciser} \add{(3.1.3)} le résultat, \struck{(3.3)}
\ill{} à l'énoncé (3.\ill{}) en écrivant $\underline{\Phi}^i_\nu = 0$ pour $i \leq n-d$.

b) $P_{<2}$ exacte, \struck{\ill{}} \uncertain{même} si $X_{\bar\eta}$ et $X_s$ sont
\struck{\uncertain{intersections complètes}} \add{\uncertain{lisses}} en dimension $n$,
\struck{et, \ill{} $n = 1$}, et si les hypothèses (3.1.1) et (3.1.2) sont satisfaites
\struck{pour} avec $m = n$, on voit \uncertain{qu'on ne peut} en général \uncertain{améliorer}
la conclusion (3.1.3) d'une nullité comme dans \uncertain{a)}.
\struck{Prenant} Prenons pour \ill{} \struck{\ill{}}, p.~ex. \struck{de $X'$ lisse}
\struck{\ill{} \uncertain{connexes} (p.~ex.} \struck{sur $S$ à fibres géom. \ill{}}
$X' = \mathbb{P}^1_S \sqcup \mathbb{P}^1_S$ et pour $X$ \uncertain{le schéma} déduit de $X'$ en
identifiant deux pts \add{distincts} (de $X'_s$ rationnels sur $k$, situés dans deux
composantes distinctes de $X'_s$. Alors \struck{les hypothèses} $X$ satisfait les hypothèses
avec $m = n = 1$, $d = 0$, cependant \struck{la conclusion} on \uncertain{n'a pas}
$\underline{\Phi}^1 = 0$, \struck{\ill{}} l'homomorphisme canonique
$H^0(X_0) \to H^0(X_{\bar\eta})$ \ill{} bijectif, puisque $X_0$ est connexe mais
$X_{\bar\eta}$ ne l'est pas.

\page{142}
\note{Pagination de l'auteur « V.15 ».}

3.3. Notre but dans la suite de ce § est de donner des résultats locaux sur la nullité
des $(\underline{\Phi}^i_\nu)$, \uncertain{tout} analogues aux résultats globaux précédents,
qui s'appliqueront \uncertain{en} \ill{} des hypothèses projectives pour $f : X \to S$. La
notion utilisée sera la \uncertain{notion} de \emph{profondeur locale} (SGA~2 XIV 4.5),
qui remplace la \uncertain{notion} de profondeur globale. Il n'est \uncertain{prouvé}
\ill{} que nous visons à \ill{} \uncertain{disposons} du th. de résolution des singularités
et du th. de \emph{profondeur affine locale}, \struck{dans divers} \uncertain{lesquels sont}
\ill{} des résultats de SGA~4 XIV, XVI, XIX, il en est \ill{} admis \struck{\ill{}})
d'\uncertain{égale} caractéristique \add{$p$}, lorsqu'on dispose de la résolution des
singularités pour les schémas \ill{} caractéristiques \struck{\ill{} \ill{}}
\add{\ill{}} \ill{} \ill{}, \uncertain{seulement} en caractéristique \emph{nulle}.
Dans le cas où \emph{$S$ est de caractéristique inégale}, \ill{} que nous disposons
\ill{} d'\uncertain{explicites} ; cela \ill{} \uncertain{à un \ill{}}
\uncertain{ingénieux} que \ill{} \ill{} \uncertain{obligé} pour le cas de la car.
\uncertain{nulle}, on en a \uncertain{déjà les} \ill{} avec des restrictions dimensionnelles
draconiennes ($\dim X_\eta \leq 2$).
\note{Les lignes centrales de ce feuillet, d'une écriture rapide, sont en grande partie
illisibles ; on n'en donne que les mots lus. Le sens général : les résultats locaux
supposent la résolution des singularités et le théorème de profondeur affine locale,
acquis en caractéristique nulle, et en caractéristique inégale seulement sous des
restrictions de dimension.}

\page{143}
\note{Pagination de l'auteur « V.16 ». Le numéro « 3 » du paragraphe est corrigé en « 4 »,
et de même dans les numéros de formules (3.x.y devenu 4.x.y) ; on donne les numéros corrigés.}

\section{4. \struck{Le cas} Situation locale des cinq schémas.}

4.1. Soient \struck{$x \in$} $x \in X_0$, $\bar x$ un pt géométrique au-dessus, $X'$ le
localisé strict de $X$ en $\bar x$, posons
\[
(4.1.1)\qquad V = X' - \{\bar x\}, \quad V_s = V \times_S s, \quad
V_\eta = V \times_S \eta = V - V_s,
\]
\[
V_{\bar\eta} = \text{\struck{$V \times_S \bar\eta$}} = V_\eta \times_\eta \bar\eta .
\]
En plus de ces quatre schémas, \struck{nous \uncertain{ferons}} nous \add{en} considérons
un cinquième. \struck{À} \add{À} cet effet, considérons le normalisé $\bar S$ de $S$ dans
$\bar K$ :
\[
(4.1.2)\qquad \bar S \longrightarrow S, \qquad \bar S = \varprojlim S_\alpha ,
\]
\struck{et posant} où les $S_\alpha$ sont les normalisés de $S$ dans les extensions finies
$K_\alpha$ de $\bar K / K$. \struck{\ill{}} On pose
\[
(4.1.3)\qquad \bar V = V \times_S \bar S = \varprojlim V_\alpha, \qquad
V_\alpha = V \times_S S_\alpha .
\]
\struck{Sont alors} Évidemment on a \struck{$\bar V_s =$} une immersion ouverte
\struck{$\bar V_\eta = \bar V \times_S \eta \simeq V_{\bar\eta} = \bar V \times_S \bar\eta \simeq V_{\bar\eta}$}
\[
(4.1.4)\qquad V_{\bar\eta} \hookrightarrow \bar V ,
\]
identifiant $V_{\bar\eta}$ à la fibre générique de $\bar V$ sur $\bar S$ [resp. sur $S$],
et le complémentaire \struck{$V_s$} $\bar V - V_{\bar\eta}$ \struck{\ill{}} s'identifie à
\struck{$V_{s}\times_s\bar s$}
\[
(4.1.5)\qquad V_{\bar s} = V_s \otimes_{k(s)} k(\bar s),
\]
où $\bar s$ est le point fermé de $\bar S$. Notons que la projection
\[
(4.1.6)\qquad V_{\bar s} \longrightarrow V_s
\]
\uncertain{induit} une \uncertain{équivalence} \uncertain{des} sites étales (SGA~4 VIII \ill{})
(car $k(\bar s)$ est une extension radicielle de $k(s)$), et \struck{permet} \add{\uncertain{permet}}
donc d'identifier $V_s$ et $V_{\bar s}$.

\drawing{Dans la marge gauche, à hauteur de (4.1.4), un croquis en cube des cinq schémas
au-dessus de $S$ : $V \leftarrow \bar V$, $V_\eta$, $V_{\bar\eta}$ en haut, $S$, $\bar S$,
$\eta$, $\bar\eta$, $s$ en bas, reliés par des flèches ; plusieurs traits sont repassés et
une partie du croquis est biffée.}

\page{144}
\note{Pagination de l'auteur « V.17 ».}

Finalement nous trouvons un diagramme commutatif \struck{à carrés cartésiens}

\begin{tikzcd}
V_s \arrow[r, hook, "i"] & V & V_\eta \arrow[l, hook', "j"'] \\
V_{\bar s} \arrow[u, "h_s"', "\wr"] \arrow[r, hook, "\bar i"] & \bar V \arrow[u, "h"] & V_{\bar\eta} \arrow[l, hook', "\bar j"'] \arrow[u, "h_\eta"']
\end{tikzcd}

(4.2), dans lequel \struck{les \uncertain{flèches} \ill{}} \uncertain{les} deux carrés
\uncertain{sont} cartésiens \ill{} au \ill{}, \ill{} qui \uncertain{concerne} $V_{\bar s}$)
$V_{\bar s}$ et $V_{\bar\eta}$ sont resp. les images inverses des \add{fermé resp. ouvert}
($V_s$ et $V_\eta$ de $V$, \struck{et} et où $h_s$ \struck{\ill{}} \add{\uncertain{en fait}}
\ill{} \uncertain{équivalence} des sites étales.

\struck{3.5. Les \uncertain{cinq}} Les cinq schémas \add{\uncertain{essentiellement} distincts}
($V_s$, $V$, $V_\eta$, $V_{\bar\eta}$, $\bar V$ jouent, du point de vue topologique, des
rôles \uncertain{bien} distincts :
\begin{itemize}
  \item $V_s$ définit la structure topologique locale de la fibre $X_0$ en $x$ ;
  \item $V$ définit la structure topologique locale du schéma ambiant $X$ en $x$ ;
  \item $V_\eta = V - V_s$ \struck{\ill{}} fournit, \add{\uncertain{\ill{}}} les
  structures cohomologiques \uncertain{reliant} $V_s$ et $V$ ;
  \item $V_{\bar\eta}$ fournit les \struck{cohomologie} \uncertain{groupes} des
  \struck{cohomologie} cycles évanescents de $X$ en $x$ [\struck{\ill{}} \add{Il est}
  relié à $V_\eta$ par $V_{\bar\eta} = V_\eta \times_\eta \bar\eta$, et à $V_s$ par la
  \uncertain{cohomologie} \ill{} la suite \struck{exacte} spectrale \uncertain{habituelle}] ;
  \item $\bar V$ permet de mettre en relation la structure cohomologique de
  $V_s \simeq V_{\bar s}$, et celle de $V_{\bar\eta}$.
\end{itemize}

\marginal{Encadré dans la marge gauche, avec une flèche vers « cycles évanescents » :
\uncertain{Noter} (4.2.1)
$(\underline{\Phi}^i_\nu)_{\bar x} \simeq H^{i-1}(V_{\bar\eta}, \Lambda_\nu)$ si $i \geq 2$,
$0 \to (\underline{\Phi}^0_\nu)_{\bar x} \to \Lambda_\nu \to H^0(V_{\bar\eta}, \Lambda_\nu)
\to (\underline{\Phi}^1_\nu)_{\bar x} \to 0$ suite exacte.}

On voit donc que les rôles de $V_\eta$ et $\bar V$

\page{145}
\note{Pagination de l'auteur « V.18 ».}

sont relativement \struck{\ill{}} \add{auxiliaires}, contrairement à \struck{celle}
\struck{Je \ill{} que} \struck{On peut dire que \ill{} schémas \ill{}} ceux de $V_s$,
$V$, $V_{\bar\eta}$. \struck{qui \ill{} \ill{} des informations} La connaissance de la
structure topologique \uncertain{de} ces derniers constitue \struck{structure topologique
des \ill{}} l'essentiel de \struck{\ill{}} l'étude topologique de $f : X \to S$ au
\uncertain{voisinage} du point $x$.

4.3. D'autre part, le groupe $I$ opère sur $\bar V$ par $S$-automorphismes,
\add{\uncertain{via} (3.4), \uncertain{via} son opération sur $\bar S$}, \uncertain{en}
induisant des opérations \emph{triviales} sur $V_{\bar s}$, et sur $V_{\bar\eta}$ les
\add{\uncertain{$V_{\bar\eta}$-}} opérations évidentes \uncertain{via}
$V_{\bar\eta} \simeq V_\eta \times_\eta \bar\eta$.

4.4. Explicitons maintenant, par des triangles exacts dans des catégories dérivées
convenables, les principales relations \struck{\ill{}} entre les cohomologies des cinq
schémas envisagés, en \uncertain{choisissant} des coefficients
$\Lambda_\nu = \mathbb{Z}/\ell^{\nu+1}\mathbb{Z}$ \uncertain{comme} \ill{} de la torsion.
Nous considérons
$\mathbb{R}\Gamma(\bar V, \Lambda_\nu)$\note{un indice $I$ est biffé sous $\Gamma$} comme élément de
$D(\Lambda_\nu[I])$, catégorie dérivée de la catégorie des $\Lambda_\nu$-\uncertain{modules}
à \uncertain{action} \uncertain{admissible} de $I$. Pour \uncertain{rappeler} ce \uncertain{fait},
nous écrivons \uncertain{plutôt} $\mathbb{R}(\Gamma, I)(\bar V, \Lambda_\nu)$ ;
l'application du foncteur \uncertain{oubli} : est déjà \uncertain{redonne}
$\mathbb{R}\Gamma(\bar V, \Lambda_\nu)$. On a alors un triangle exact

\begin{tikzcd}[column sep=small]
 & \mathbb{R}(\Gamma, I)(V_{\bar\eta}, \Lambda_\nu) \arrow[dl] & \\
\mathbb{R}(\Gamma, I)(V_{\bar s}, \underline{\Phi}^\bullet_{V_s}) \arrow[rr] & & \mathbb{R}(\Gamma, I)(\bar V, \Lambda_\nu) \arrow[ul, "\bar j^*"']
\end{tikzcd}

(A), où $\underline{\Phi}^\bullet_{V_s} = \mathbb{R}\Gamma_{V_{\bar s}}(\Lambda_{\nu\,\bar V})$,
peut aussi s'interpréter comme
\note{Le numéro (A) est écrit dans la marge, sur un numéro biffé illisible. La flèche
descendante gauche porte une petite marque, sans doute le degré $+1$ du triangle.}

\page{146}
\note{Pagination de l'auteur « V.19 ».}

la restriction à $V_s$ du complexe $\underline{\Phi}^\bullet_\nu$ \add{à opérateurs $I$}
sur $X_0$, déjà \add{\uncertain{introduit}} \uncertain{introduit} dans III~1 (du moins,
pour ses faisceaux de cohomologie, qui y \struck{\ill{}} \add{sont} notés
$\underline{\Phi}^q_\nu$). Par application du foncteur oubli, on \uncertain{trouve} un triangle
\uncertain{dérivé}

\begin{tikzcd}[column sep=small]
 & \mathbb{R}\Gamma(V_{\bar\eta}, \Lambda_\nu) \arrow[dl] & \\
\mathbb{R}\Gamma(V_s, \underline{\Phi}^\bullet_\nu) \arrow[rr] & & \mathbb{R}\Gamma(\bar V, \Lambda_\nu) \arrow[ul, "\bar j^*"']
\end{tikzcd}

(A$'$). \struck{(A) le foncteur $\mathbb{R}\Gamma^I$, \ill{}} \struck{Lorsque}
\struck{triangle exact} Supposons que
\[
(4.4.1)\qquad (\underline{\Phi}^i_\nu)_{\bar x} = 0 \quad \text{pour } i < m ,
\]
\struck{la suite exacte \ill{}} \uncertain{alors} (A$'$) donne
\[
(4.4.2)\qquad H^i(V_s, \underline{\Phi}^\bullet_\nu) = 0 \quad \text{pour } i < m ,
\]
et la suite exacte : (A$'$) donne
\[
(4.4.3)\qquad H^i(\bar V, \Lambda_\nu) \xrightarrow{\ \bar j^*\ } H^i(V_{\bar\eta}, \Lambda_\nu)
\quad \begin{cases} \text{injectif pour } i \leq m-1 \\ \text{bijectif pour } i \leq m-2 . \end{cases}
\]
\note{Les deux conditions sont entourées ; l'inégalité de la première ligne est écrite
sur une autre biffée, et pourrait être $i = m-1$.}

Si (4.4.1) est vrai pour $m = +\infty$, i.e. si $x$ \uncertain{est un} pt \struck{\ill{}}
de l'ens. des \uncertain{points} de \uncertain{locale acyclicité} de $f : X \to S$, on trouve
donc
\[
\mathbb{R}(\Gamma, I)(V_s, \underline{\Phi}^\bullet_\nu) = 0 ,
\]
i.e. que \struck{$\mathbb{R}\Gamma, I$}
\[
\mathbb{R}(\Gamma, I)(\bar V, \Lambda_\nu) \xrightarrow{\ \bar j^*\ }
\mathbb{R}(\Gamma, I)(V_{\bar\eta}, \Lambda_\nu)
\]
est un \uncertain{isom.}, i.e.
\[
H^i(\bar V, \Lambda_\nu) \xrightarrow[\sim]{\ \bar j^*\ } H^i(V_{\bar\eta}, \Lambda_\nu)
\quad \text{pour tout } i ,
\]
\uncertain{on \ill{}} (4.4.\ill{}) \ill{} \uncertain{pour} tout $i$.

\marginal{Encadré dans la marge gauche, écrit le long du bord, sous le titre (A$'$) :
\ill{} qui \uncertain{peut} aussi s'écrire \ill{} (4.4.4)
$H^i_x(\bar V, \Lambda_\nu) \to (\underline{\Phi}^i_\nu)_{\bar x}$ bijectif si $i < m$ \ill{},
compte tenu de (4.4.5)
$H^i_{\bar x}(\bar X, \Lambda_\nu) \simeq \varinjlim_m H^i_x(V(m), \Lambda_\nu)$
\ill{}, $H^i(\bar V, \Lambda_\nu) \to H^i(V_{\bar\eta}, \Lambda_\nu)$ déduite de la relation
4.2.1 et des \uncertain{relations} \ill{}, déduite des \uncertain{homomorphismes}
$V_{\bar\eta} \to \bar V$ \ill{}.}
\note{L'encadré marginal, écrit perpendiculairement au texte et en partie recouvert par
une accolade, n'est lu que par fragments.}

\page{147}
\note{Pagination de l'auteur « V.20 ».}

Dans ce dernier cas, \struck{\ill{}} l'inclusion $V_{\bar\eta} \to \bar V$ induit un
isomorphisme sur \uncertain{le} \uncertain{type} cohomologique (à coeff. dans $\Lambda_\nu$),
et on peut dire que les cinq schémas initiaux \struck{\ill{}} se réduisent à quatre,
$V_s$, $V$, $V_\eta$, $V_{\bar\eta}$ (on est débarrassé du seul $\bar V$ des cinq schémas
qui n'est pas \uncertain{noethérien} !).

4.5. L'inclusion $V_{\bar s} \hookrightarrow \bar V$ définit de même un homomorphisme
\[
(4.5.1)\qquad \mathbb{R}(\Gamma, I)(\bar V, \Lambda_\nu) \longrightarrow
\mathbb{R}(\Gamma, I)(V_s, \Lambda_\nu)
\]
qui s'insère dans un triangle exact

\begin{tikzcd}[column sep=tiny, nodes={font=\small}]
 & \mathbb{R}(\Gamma, I)(V_{\bar s}, \Lambda_\nu) \simeq \mathbb{R}\Gamma(V_s, \Lambda_\nu)^{\mathrm{triv}} \arrow[dl] & \\
\mathbb{R}(\Gamma, I)(\bar V, (\Lambda_\nu)_{V_{\bar\eta}, \bar V}) \arrow[rr] & & \mathbb{R}(\Gamma, I)(\bar V, \Lambda_\nu) \arrow[ul, "\bar i^*"']
\end{tikzcd}

(B).\note{Le numéro (B) remplace dans la marge un numéro biffé, peut-être « (4.2.3) ».
Dans le terme du haut, $V_{\bar s}$ est écrit sur un autre indice biffé.}

L'application du foncteur oubli $I$ \uncertain{au} diagramme donne

\begin{tikzcd}[column sep=small]
 & \mathbb{R}\Gamma(V_s, \Lambda_\nu) \arrow[dl] & \\
\mathbb{R}\Gamma(\bar V, (\Lambda_\nu)_{V_{\bar\eta}, \bar V}) \arrow[rr] & & \mathbb{R}\Gamma(\bar V, \Lambda_\nu) \arrow[ul, "\bar i^*"']
\end{tikzcd}

(B$'$).

\struck{Sous la \uncertain{condition} (4.5.2)}

\page{148}
\note{Pagination de l'auteur « V.21 ».}

\section{5. Application du th. de profondeur locale (SGA~2 XIV 4).}

5.0. Cf. les commentaires dans 3.3. sur les hypothèses faites
\add{quand nous appliquerons les théorèmes de profondeur locaux}.

5.1. Soit $f : X \to S$ loc. de type fini, $S$ str. local hensélien, et soit
$x \in X_0$. Soit $m' \in \mathbb{Z}$ et supposons que l'on ait
\struck{(\uncertain{conditions} d'hypothèse}
\[
(5.1.1)\qquad (\underline{\Phi}^i_\nu)_{\bar x'} = 0
\]
pour \struck{\ill{}} tout pt géom. $\bar x'$ \uncertain{voisin} de $x$ dans $X_0$, et
$i \leq m'$. \struck{On} Nous considérons cette hypothèse comme une hypothèse de
\uncertain{récurrence}, et nous proposons de donner des conditions moyennant \ill{}
\uncertain{conclure}
\[
(*)\qquad (\underline{\Phi}^i_\nu)_{\bar x} = 0 \quad \text{si } i \leq m .
\]
Notons que l'on sait déjà que (5.1.1) implique
\[
\text{\struck{$(5.1.1)\quad H^i(\bar V, \Lambda_\nu) \xrightarrow{\sim} H^i(X_{\bar\eta}, \Lambda_\nu)$ si $i < m'$}}
\]
\struck{\ill{}}
\[
(5.1.2)\qquad (\underline{\Phi}^i_\nu)_{\bar x} = H^i_x(\bar X', \Lambda_\nu)
= \varinjlim_h H^i_x(X(h), \Lambda_\nu) \quad \text{pour } i \leq m' .
\]
Par suite, \struck{\ill{}} la relation $(*)$ \struck{\uncertain{équivaut}} \uncertain{résulterait}
de la relation $H^i_x(X(h), \Lambda_\nu) = 0$ pour $i \leq m'$, $\forall h$, i.e.
\struck{(5.1.3)} de la relation
\[
(5.1.3)\qquad \operatorname{prof}_x(X(h)) \geq m'+1 \quad \text{pour tout } h .
\]
\note{Les numéros (5.1.1) et (5.1.3) sont cerclés dans la marge.}

D'où \struck{\ill{}} on conclut aussitôt, \add{par récurrence sur
$\dim \mathcal{O}_{X_0,x} = r$ :} \uncertain{l'énoncé suivant}.

\emph{\struck{Proposition} Théorème} 5.2. Soit $Z$ une partie de $X_0$,
(\uncertain{supposons} \ill{}. Supposons que :

a) Pour toute générisation $x'$ de $x$ dans $X_0$, \struck{distincte de $x$,}
\add{contenue} dans $X_0 - Z$, on a $(\underline{\Phi}^i_\nu)_{\bar x'} = 0$ si $i \leq m'$.

\page{149}
\note{Pagination de l'auteur « V.22 ».}

(c'est le cas si on suppose p.~ex. que $Z$ est l'ensemble des pts de $X_0$ en lesquels
$f$ n'est pas lisse).

b) Pour toute générisation $x'$ de $x$ dans $Z$, \struck{distincte} on a
\struck{\ill{}} \add{\uncertain{et tout}} $h$,
\[
\operatorname{prof}_{\bar x'}(X(h)) \geq m'+1 ,
\]
\struck{(ce qui est le cas} Alors on a
\[
(5.2.1)\qquad (\underline{\Phi}^i_\nu)_x = 0 \quad \text{si } i \leq m' ,
\]
i.e. $f$ est loc. $(m'+1)$-acyclique en $x$, pour \struck{\ill{}} $\{\ell\}$, donc si $f$ est
propre, on a
\[
(5.2.2)\qquad H^i(X_0, \Lambda_\nu) \longrightarrow H^i(X_{\bar\eta}, \Lambda_\nu)
\]
un isom. si $i < m'$, mon. si $i = m'$.

\note{Toute la suite du feuillet, depuis « Corollaire 5.3 » jusqu'en bas, est encadrée et
biffée de deux longues diagonales ; on la donne entière comme passage biffé.}
\struck{\emph{Corollaire} 5.3. Soit $Z$ une partie de $X_0$, \ill{}, supposons}
\struck{$d = \sup \dim \overline{\{x'\}}$ ($x' \in Z$, $x'$ générisation de $x$),}
\struck{(5.3.1) $m = \ill{}$, $1+m = \sup_h \ill{}$.}
\struck{Alors on a, supposons que a) $(\underline{\Phi}^i_\nu)_{\bar x'} = 0$ si
$i \leq m-d$, pour $x' \in X - Z$, $x \in \overline{\{x'\}}$,}
\struck{b) Pour toute générisation $x'$ de $x$ dans $X_0$, \ill{}
$\operatorname{prof}^x_{x'}(X_0) \geq m+1$ si $x' \in X_0$ \uncertain{\ill{}}
$\operatorname{prof}^x_{x'}(X_y) = 0$ si $x' \in X_y$).}
\struck{i.e. $f$ est loc. $(m-d)$-acyclique en $x$, pour $\{\ell\}$.}
\struck{Alors}
\struck{Résulte formellement de 5.2, en y faisant $m' = m-d$, la condition a) de 5.2
étant satisfaite par \uncertain{\ill{}} de (5.3.2), et la condition b) résultant de}

\page{150}
\note{Pagination de l'auteur « V.23 ».}

\marginal{Encadré en haut à gauche : \struck{Supposons} la condition a) de 5.2 satisfaite.}
(Cf. 5.0. sur les hypothèses préliminaires) \struck{\ill{}} \add{de 5.2.}
(dans la conclusion 5.2.1)

\emph{Corollaire} 5.3. La condition b) est satisfaite dans chacun des cas suivants (i) et
(ii), où on a posé
$d = \sup \dim \overline{\{x'\}}$, \struck{$x' \in Z$}, où $x'$ parcourt les
générisations de $x$ dans $X_0$ :

\struck{(i)} \emph{Cas I}. Pour toute générisation $x'$ de $x$ dans $X$, on a
\[
(5.3.1)\qquad
\begin{cases}
\operatorname{prof\,et}^*_{x'}(X_0) \geq m'+1 & \text{si } x' \in Z \\
\operatorname{prof\,et}^*_{x'}(X_{\bar\eta}) \geq m' + \dim \overline{\{x\}} & \text{si } x' \in X_{\bar\eta} ,
\end{cases}
\]
ce qui est le cas si on \uncertain{suppose} \uncertain{que} $m' = m-d-1$, avec $m$
un entier tel que
\[
(5.3.2)\qquad
\begin{cases}
\operatorname{prof\,et}^x_{x'}(X_0) \geq m & \text{pour toute \add{générisation} } x' \in Z \text{ de } x \\
\operatorname{prof\,et}^x_{x'}(X_{\bar\eta}) \geq m-1 & \text{pour toute gén. $x'$ de $x$ dans $X$, avec } x' \in X_{\bar\eta} .
\end{cases}
\]
\note{L'astérisque en exposant de « prof et » est écrit sur un autre signe ; dans
(5.3.1) le second membre de la deuxième ligne porte une rature avant « $+\dim$ ».}

\emph{Cas II}. \struck{Pour toute \uncertain{générisation}} \struck{On a} \add{Le morphisme
$f$ est plat en $x$, et on a} \struck{\ill{} générisation $x'$ de $x$ dans $Z$,}
\[
\text{\struck{$\ill{}$}}\quad m' = \operatorname{prof\,géom}^*_x(X_0) - d ,
\]
où $\operatorname{prof\,géom}^*_x(X_0) \overset{\mathrm{df}}{=}
\operatorname{prof\,géom}_x(X_0) + \dim \overline{\{x\}}$
(cf. SGA~2 XIV 5.3 pour la notion de profondeur géométrique),
\struck{\ill{} le cas si on fait \ill{}, \ill{} entier \ill{} tel que
$\operatorname{prof\,géom}^*_x(X)$}.
\note{Un passage encadré, à la fin du cas II, est biffé de traits croisés ; on n'en lit
que « prof géom$^*$ ».}

\emph{Démonstration}. \emph{Cas I.} Les relations (5.3.1) \struck{et} \add{\uncertain{et}} se
\uncertain{stabilisent} \add{par générisation et} par changement de $X$ en $X(h)$, il suffit de
vérifier qu'elles \struck{\ill{}} entraînent
\[
(*)\qquad \operatorname{prof\,et}_x(X) \geq m'+1 .
\]
Or la deuxième des relations (5.3.1) \add{s'écrit} \struck{\ill{}}
\struck{s'écrivent}
\[
\operatorname{prof\,et}_{x'}(X') + \dim \overline{\{x'\}}^{(X_{\bar\eta})}
\geq m' + \dim \overline{\{x\}} ,
\]
or
\[
\dim \overline{\{x'\}}^{(X_{\bar\eta})} = \dim \overline{\{x'\}}^{X} - 1
= \dim \overline{\{x\}} + \dim \overline{\{\bar x'\}}^{(X')} - 1
\]
où $\bar x'$ est un pt de $X'$ au-dessus de $x'$, donc on \uncertain{trouve}
\note{Au-dessus de la première égalité, en marge droite, une variante :
« $= \dim \overline{\{x\}} + \dim \overline{\{\bar x'\}}^{(V)}$ ».}

\page{151}
\note{Pagination de l'auteur « V.24 ».}

l'interpréter en disant que pour $x' \in X'_{\bar\eta} = X' - X'_0$, on a
\[
\operatorname{prof\,et}_{x'}(V) + \dim \overline{\{x'\}}^{(V)} \geq m'
\ \text{\struck{$\ill{}$}} ,
\]
ce qui par le th. de profondeur locale SGA~1 XIV 4.5
\note{sic, sans doute pour SGA~2.}
implique que
\[
H^i(V) \longrightarrow H^i(V_s) \quad \text{est}\quad
\begin{cases} \text{bijectif pour } i < m'-1 \\ \text{injectif pour } i = m'-1 . \end{cases}
\]
\marginal{Dans la marge gauche, en oblique : il suffirait que ce soit bijectif pour
$i \leq m'-1$ \ill{}, injectif pour $i = m'$ \ill{}.}
Donc il \struck{suffit, pour la relation \ill{}} \uncertain{est} \struck{\ill{} que}
\struck{\ill{}} équivalent à
$\operatorname{prof\,et}^*_x(X_0) \geq m'+1$, ce qui est la première des relations (5.3.1).
Le fait que les relations, (5.3.2) impliquent les relations (5.3.1) est d'autre part
immédiat.

\emph{Cas II.} Dans ce cas, l'hypothèse est \add{évidemment} stable par passage de $X$
à $X(h)$, \struck{\ill{}} il suffit de \add{\ill{}} \uncertain{vérifier} \struck{\ill{}}
\struck{\ill{} \uncertain{générisation}, \ill{} \uncertain{condition} \ill{}}
\uncertain{qu'elle} implique \ill{}
\[
(**)\qquad \text{\struck{$\operatorname{prof\,et}_{x'}(X) \geq m'+1$}}
\]
pour toute génér. $x'$ de $x$ dans $Z$. Or on a (\uncertain{lemme}
\struck{\ill{}} 5.3.3 (ii) (plus bas !))
\[
\text{\struck{$\operatorname{prof\,et}_{x'}(X) \geq \operatorname{prof\,et}'_x$}}
\]
\[
\operatorname{prof\,géom}_{x'}(X_0) \geq \operatorname{prof\,géom}_x(X)
\ \text{\struck{\ill{}}}\ \operatorname{codim}(\overline{\{x\}}, \overline{\{x'\}})
\]
et, comme
$\operatorname{codim}(\overline{\{x\}}, \overline{\{x'\}}) = \dim \overline{\{x'\}} - \dim \overline{\{x\}}
\leq d - \dim \overline{\{x\}}$, on trouve
\[
\operatorname{prof\,géom}_{x'}(X_0) \geq \operatorname{prof\,géom}^*_x(X_0)
\ \text{\struck{\ill{}}}\ - d = m' .
\]
D'autre part, \uncertain{par le lemme} 5.3.3 (ii) (plus bas), on \ill{}
\[
\operatorname{prof\,géom}_{x'}(X) \geq \operatorname{prof\,géom}_{x'}(X_0) + 1 ,
\]
\note{Sous « $+1$ », une glose : « $= \operatorname{prof\,géom}_s(S)$ ».}
donc on en conclut
\[
\operatorname{prof\,géom}_{x'}(X) \geq m'+1 .
\]

\page{152}
\note{Pagination de l'auteur « V.25 ».}

D'autre part, le th. de profondeur locale entraîne (SGA~1 XIV 5.6)
\note{sic, sans doute pour SGA~2.}
\[
\operatorname{prof\,et}_{x'}(X) \geq \operatorname{prof\,géom}_{x'}(X) ,
\]
d'où \struck{\ill{}} pour $(**)$, cqfd.

Je vérifie : conjectures suivantes \ill{}

\emph{Lemme} 5.3.3. (i) Soit $A$ un anneau local noethérien, $\mathfrak{p}$ un idéal premier
de $A$. Alors on a
\[
\operatorname{prof\,géom} A_{\mathfrak{p}} \geq \operatorname{prof\,géom} A - \dim A/\mathfrak{p} .
\]

(ii) Soit $A \to B$ un \ill{} local d'anneaux locaux noethériens, avec $B$
\emph{plat} sur $A$, $k$ le corps résiduel de $A$, $K$ celui de $B$, et supposons qu'il
existe un anneau local noethérien $C$ \struck{\ill{} $k$, \ill{}}
\uncertain{formellement} \struck{lisse} sur $k$ de corps résiduel $K$ (\uncertain{condition}
\uncertain{satisfaite} \uncertain{toujours} si $K$ est une extension de type fini de $k$).
Alors on a
\[
\operatorname{prof\,géom} B \geq \operatorname{prof\,géom} A + \operatorname{prof\,géom} B \otimes_A k .
\]

Donner démonstration, !

\note{Le reste du feuillet est blanc, sauf une note en bas de page.}
\marginal{En bas de page : $(*)$ Cette condition est peut-être inutile. Pour une
\uncertain{généralisation} cf. EGA~0$_{\mathrm{IV}}$ 22.2.2.}

\page{153}
\note{Pagination de l'auteur « V.26 ».}

Voici un corollaire du cas II :

\emph{Corollaire} 5.4. Supposons \add{$X$ plat sur $S$ et} $X_0$ loc. intersection complète
et purement de dimension $n$. Alors on a
\[
\underline{\Phi}^i_\nu = 0 \quad \text{pour}\quad
i \notin [n-d+1, n+d+1] ,
\]
\note{Au-dessus, une première version biffée :
« $i \notin [n-d+1, n+1]$ » puis « $0 \leq n-d$ » ; une variante
« $i \notin [n-d+1, n+1]$ » est récrite au-dessus de la ligne biffée.}
où $d$ est la dimension de l'ens. $Z$ des pts de $X_0$ en lesquels $f$ n'est pas lisse,
\struck{\ill{}} et
\[
\underline{\Phi}^i_\nu = 0 \quad \text{si } i \neq n+1 .
\]
En particulier, si $d = 0$, on a
\[
\underline{\Phi}^i_\nu = 0 \quad \text{si } i \neq n+1 ,
\]
donc l'homomorphisme
\[
H^i(X_0, \Lambda_\nu) \longrightarrow H^i(X_{\bar\eta}, \Lambda_\nu)
\]
est un isomorphisme si $i \neq n, n+1$, un monomorphisme pour $i = n$, \add{de plus}
un épimorphisme pour $i = n+1$ ; (\uncertain{\ill{}} cela, \uncertain{dans} \ill{}
pour \uncertain{dimension} $I$ \uncertain{sur} $H^i(X_{\bar\eta}, \Lambda_\nu)$
\struck{est \uncertain{triviale}}, \uncertain{\ill{}} $H^i(X_{\bar\eta}, \mathbb{Z}_\ell)$)
est triviale si $i \neq n$.
\note{La parenthèse finale, serrée entre les lignes, n'est lue que par fragments.}

5.5. Pour terminer ce paragraphe, nous allons donner une variante des résultats précédents
pour la cohomologie non commutative. Supposons \struck{\ill{}} \add{pour fixer les idées}
$f$ loc. 1-acyclique en les générisations \add{distinctes de $x$} de $x$, \struck{\ill{}}
\ill{} trouver \struck{\ill{}} \uncertain{conditions} \struck{\ill{}} de
\uncertain{locale} 1-asphéricité en $x$, pour \struck{\ill{}} \uncertain{l'un} et
\uncertain{pour} \uncertain{prendre} \struck{\ill{}} \add{$\mathbb{L}$}.
\struck{\ill{}}
Je dis que pour tout $\mathbb{L}$-\add{groupe} fini $G$, \struck{\ill{}} \ill{} :

\begin{tikzcd}
V \arrow[r, no head] \arrow[d, no head, "\cap" description] & \bar V \arrow[dd, no head] \\
X' \arrow[d, no head] & \\
S & \bar S \arrow[l]
\end{tikzcd}

\[
(5.5.1)\qquad H^i(\bar V, G) \xrightarrow{\ \sim\ } H^i(V_{\bar\eta}, G)
\quad \text{si } i \leq 1 .
\]
\note{Le petit diagramme est dans la marge gauche, en face de 5.5 ; seule la flèche
$\bar S \to S$ porte une pointe, les autres traits sont transcrits sans tête, et le signe
entre $V$ et $X'$ est lu comme une inclusion.}

\page{154}
\note{Pagination de l'auteur « V.27 ».}

Donc pour avoir la locale 1-asphéricité pour $\mathbb{L}$ en $x$, il suffit de supposer
que \uncertain{l'on a}, pour tout tel $G$ :
\[
(5.5.2)\qquad H^i(V(h), G) = 0 \quad \text{si } i \leq 1 .
\]
On sait \struck{\ill{}} que cette condition est vérifiée \struck{\ill{}} \struck{les}
\struck{$X(h)$ sont} \add{les} intersections complètes de dim $\geq 3$ en $x$,
(c'est sans doute dire que leur prof géom en $x$ est $\geq 3$). Lorsque $f$ est plat en $x$,
cela signifie aussi que $X_0$ est une int. complète de dimension $\geq 2$ en $x$.
Utilisant le th. de locale 1-asphéricité pour les morphismes lisses, on trouve ainsi le

\emph{Théorème} 5.6. Supposons que $f : X \to S$ soit plat, \uncertain{pour} l'ens. $Z$
des points \add{de $X_0$} en lesquels $f$ n'est pas lisse, \uncertain{soit} de codimension
$\geq 2$, et que $X_0$ soit une intersection complète locale en tout ses points. Alors $f$
est localement 1-asphérique pour \add{$\mathbb{L}$} \struck{\ill{}}
$\mathbb{L} \neq \mathbb{P} - \{p\}$. En particulier si $f$ est propre, \uncertain{alors}
pour tout groupe fini $G$ (\add{d'ordre} premier à $p$), l'homomorphisme de
spécialisation
\[
H^1(X_0, G) \longrightarrow H^1(X_{\bar\eta}, G)
\]
est un isomorphisme.

\emph{Remarque} 5.7. a) Je doute qu'il doive être possible de remplacer l'hypothèse
que $X_0$ soit une intersection complète par celle que pour tout $x \in Z$, on ait
$\operatorname{prof\,géom}_x(X_0) \geq 2$, cela résultant d'une relation
\uncertain{affirmative} : SGA~1 XIV 6.3
\note{Le feuillet s'achève sur ce deux-points ; la suite au feuillet suivant.}

\page{155}
\note{Pagination de l'auteur « V.28 ». Suite de la remarque 5.7.}

b) Les résultats précédents \uncertain{suggèrent} la conjecture suivante
\add{(au moins pour $A$, $B$ \emph{excellents})} :

Soit $f : X \to S$ un morphisme \emph{plat} de schémas
\add{$= \operatorname{Spec} B$, $= \operatorname{Spec} A$} \add{strictement locaux},
noethériens, \struck{\ill{} à fibres géométriques \ill{}} $s$ le point fermé de $S$,
$X_0 = f^{-1}(s)$, \struck{\ill{} $X_0$, $k$ le corps, \ill{} le pt \ill{}}
\struck{condition $(R_d)$} \add{\ill{}} \struck{fait} $(R_d)$, et que pour tout
\struck{$X_0$ soit géométriquement \ill{}} extension finie $k'$ de $k = k(s)$, l'ens.
des points singuliers de $X_0 \otimes_k k'$ soit de dim $\leq d$,
\struck{\ill{} condition $R_d$, et \uncertain{supposons}} posons
$m = \operatorname{prof\,géom} B$. Alors $f$ est loc. $(m-d-1)$-acyclique,
\struck{et si $m-d-1 \geq 1$} pour $\mathbb{L} = \mathbb{P} - \{p\}$, $p$ car. résiduelle
de $A$), et si $m-d-1 \geq 1$ i.e. $m-d \geq 2$, $f$ est loc. 1-asphérique pour
$\mathbb{L}$. \struck{Cela} Il est plausible que cela soit démontrable comme conséquence
de la conjecture déjà citée SGA~1 XIV 6.3, et de la résolution des singularités pour les
schémas finis sur $S$.
\note{Les lignes du milieu, surchargées de ratures, ne sont lues que partiellement ;
l'ordre des clauses suit la lecture la plus probable.}

\page{156}
\note{Pagination de l'auteur « V.29 ». À partir d'ici le groupe opérant est noté $\pi$ ; p.~157 il précise $\pi = I$.}

\section{6. Cas d'un point \struck{\ill{}} isolé \add{dans $X_0$} non lisse maximal \struck{\ill{}}}

(6.1) Soit $x \in X_0$, et \struck{supposons} reprenons les notations du §4.
Nous supposons que l'on a pour tout $i$
\[
(6.1.1)\qquad \underline{\Phi}^i_\nu | V = 0, \ \text{i.e.}\ 
(\underline{\Phi}^i_\nu)_{\bar x'} = 0 \ \text{pour toute gén. } x' \neq x \text{ dans } X_0 .
\]
Condition satisfaite en particulier si $x$ est un point maximal de l'ensemble
\struck{$Z$} des pts de $X_0$ en lesquels $f$ n'est pas lisse,
\struck{\uncertain{et} \uncertain{lorsque} $f$ \ill{} plat, il \ill{} (4.4.3) \ill{}
\ill{} \ill{} de $X$ dominant \ill{} et \ill{} $x$}.
\note{Un passage encadré et biffé, renvoyant à (4.4.3), n'est pas lu.}
L'inclusion
\[
V_{\bar\eta} \longleftarrow \bar V
\]
définit alors un isomorphisme
\[
(6.1.2)\qquad \mathbb{R}\Gamma(\bar V, \Lambda_\nu) \xrightarrow{\ \sim\ }
\mathbb{R}\Gamma(V_{\bar\eta}, \Lambda_\nu)
\]
i.e. des isomorphismes
\[
(6.1.3)\qquad H^i(\bar V, \Lambda_\nu) \simeq H^i(V_{\bar\eta}, \Lambda_\nu) .
\]
\note{La flèche « $V_{\bar\eta} \leftarrow \bar V$ » est orientée ainsi sur la page,
à rebours de l'inclusion (4.1.4).}
On peut en un certain \uncertain{sens} \uncertain{préciser} $\mathbb{R}$, (6.1.2)
\struck{\ill{}} s'écrit
\[
(6.1.4)\qquad \mathbb{R}(\Gamma, \pi)(\bar V, \Lambda_\nu) \xrightarrow{\ \sim\ }
\mathbb{R}(\Gamma, \pi)(V_{\bar\eta}, \Lambda_\nu) .
\]
Donc, au point de vue de leur structure cohomologique, la situation des 5 schémas (4.2)
se réduit aux \uncertain{quatre} schémas $V_s \simeq V_{\bar s}$, $V$, $V_\eta$, $V_{\bar\eta}$.

\emph{(6.2)} L'inclusion $V_{\bar s} \to \bar V$ donne donc un homomorphisme
\[
(6.2.1)\qquad u = (k^*_\pi)^* : \mathbb{R}(\Gamma, \pi)(V_{\bar\eta}, \Lambda_\nu)
\longrightarrow \mathbb{R}(\Gamma, \pi)(V_{\bar s}, \Lambda_\nu) ,
\]
\note{La flèche (6.2.1) est encadrée ; au-dessus, deux gloses entourées identifient la
source à $K^\bullet$ et le but à « $(L^\bullet)^{\mathrm{triv}}$, où
$L^\bullet = \mathbb{R}\Gamma(V_s, \Lambda_\nu)$ ».}
qui \struck{\ill{}} s'insère dans le triangle distingué

\begin{tikzcd}[column sep=tiny, nodes={font=\small}]
 & \mathbb{R}(\Gamma, \pi)(V_{\bar s}, \Lambda_\nu) \simeq \mathbb{R}\Gamma(V_s, \Lambda_\nu)^{\mathrm{triv}} = L^{\bullet\,\mathrm{triv}} \arrow[dl] & \\
R^\bullet = \mathbb{R}(\Gamma, \pi)(\bar V, (\Lambda_\nu)_{V_{\bar\eta}, \bar V}) \arrow[rr] & & \mathbb{R}(\Gamma, \pi)(V_{\bar\eta}, \Lambda_\nu) = K^\bullet \arrow[ul]
\end{tikzcd}

(6.2.2).

\marginal{Encadré en bas à gauche : compte tenu que $\pi$ opère trivialement sur
$V_{\bar s} = V_s$.}

\page{157}
\note{Pagination de l'auteur « V.30 ».}

L'homomorphisme $u$, (\add{ou ce qui revient au même,} \struck{\ill{}} le triangle
(6.2.2) correspondant) domine la structure cohomologique des quatre schémas
$V_s \simeq V_{\bar s}$, $V$, $V_\eta$, $V_{\bar\eta}$, en ce sens qu'il permet de
retrouver formellement les cohomologies (à coefficients constants) des schémas
envisagés. Appliquant le foncteur « oubli des $\pi$ » à 6.2.2, on trouve le triangle exact

\begin{tikzcd}[column sep=small]
 & \mathbb{R}\Gamma(V_s, \Lambda_\nu) = L^\bullet \arrow[dl] & \\
R^{\bullet\,\mathrm{ab}} = \mathbb{R}\Gamma(\bar V, (\Lambda_\nu)_{V_{\bar\eta}, \bar V}) \arrow[rr] & & \mathbb{R}\Gamma(V_{\bar\eta}, \Lambda_\nu) = K^{\bullet\,\mathrm{ab}} \arrow[ul, "k^*"']
\end{tikzcd}

(6.2.3), qui donne la cohomologie (sans opérations) de $V_{\bar\eta}$ et $V_s$ et
l'homomorphisme du premier dans le second défini par l'inclusion
\[
V_{\bar s} \xrightarrow{\ \bar i\ } \bar V \overset{\mathrm{coh.}}{\simeq} V_{\bar\eta} ,
\]
\note{Sous cette ligne, une flèche pointillée notée $k$ va de $V_{\bar s}$ à $V_{\bar\eta}$ :
la « flèche » cohomologique obtenue en composant l'inclusion avec l'inverse de
l'isomorphisme cohomologique (6.1.2).}
D'autre part, appliquant à (6.2.1) \add{ou (6.2.2)} le foncteur $\mathbb{R}\Gamma^\pi$,
on trouve

\begin{tikzcd}[column sep=tiny, nodes={font=\small}]
 & \mathbb{R}\Gamma(V_s, \Lambda_\nu) \overset{\mathbb{L}}{\otimes}_{\Lambda_\nu} \mathbb{R}\Gamma^\pi(\Lambda_\nu) \simeq L^\bullet \oplus L^\bullet[-1] \arrow[dl] & \\
\mathbb{R}\Gamma^\pi(R^\bullet) \arrow[rr] & & \mathbb{R}\Gamma(V_\eta, \Lambda_\nu) \simeq \mathbb{R}\Gamma^\pi(K^\bullet) \arrow[ul, "\mathbb{R}\Gamma^\pi k^*"']
\end{tikzcd}

(6.2.4), compte tenu \struck{\ill{}} que, par la \uncertain{dimension} \uncertain{cohomologique}
de $\pi = I$, \uncertain{on a}
\[
(6.2.5)\qquad \mathbb{R}\Gamma^\pi(\Lambda_\nu) \simeq \Lambda_\nu \oplus \Lambda_\nu(-1)[-1] .
\]
\note{Gloses de sa main : sous le dernier terme de (6.2.4), « dégénérescence de Leray » ;
au-dessus de « $[-1]$ » de (6.2.5), « translat. des degrés » ; sous « $\Lambda_\nu(-1)$ »,
« twist de Tate ». Dans le terme du haut de (6.2.4), un symbole biffé précède
$L^\bullet \oplus L^\bullet[-1]$.}

\struck{L'homomorphisme $\mathbb{R}\Gamma^\pi(k^\#)$ \uncertain{s'identifie} \uncertain{à}
\uncertain{\ill{}} un $\mathbb{R}\Gamma(V_\eta, \Lambda_\nu) \longrightarrow$}
Ainsi on retrouve $\mathbb{R}\Gamma(V_\eta, \Lambda_\nu)$, et l'homomorphisme $h^*$
\uncertain{déduit} de $V_{\bar\eta} \xrightarrow{h} V_\eta$ n'est autre que l'homomorphisme
\uncertain{naturel} « d'inclusion » $\mathbb{R}\Gamma^\pi(K^\bullet) \to K^\bullet$.
De plus, l'homomorphisme $\mathbb{R}\Gamma^\pi(k^*)$ se décompose en
\[
(6.2.6)\qquad \mathbb{R}\Gamma^\pi(k^*) = (\alpha, \beta) :
\mathbb{R}\Gamma(V_\eta, \Lambda_\nu) = \mathbb{R}\Gamma^\pi(K^\bullet)
\]
\[
\longrightarrow \mathbb{R}\Gamma(V_s, \Lambda_\nu) + \mathbb{R}\Gamma(V_s, \Lambda_\nu(-1))
\]
\note{La formule (6.2.6) continue au feuillet suivant.}

\drawing{Encadré dans la marge gauche, avec l'en-tête « remarquer le diagramme et le
\uncertain{commutatif} ! » écrit en oblique : un carré
$V \leftarrow V_\eta$ (flèche épaisse, étiquette biffée), $V_s \simeq V_{\bar s} \to V$
(notée $i$), $\bar V \simeq V_{\bar\eta} \to V_\eta$ (notée $h_\eta$), et en bas une flèche
pointillée $k$ de $V_s \simeq V_{\bar s}$ vers $\bar V \simeq V_{\bar\eta}$. Sous le carré :
« (avec des notations légèrement différentes de celles de (4.2)) ».}

\page{158}
\note{Pagination de l'auteur « V.31 ». Suite de (6.2.6).}

où $\alpha$ \struck{\ill{}} est donné par
\[
(6.2.7)\qquad \alpha = k^* h^* ,
\]
tandis que $\beta$ est un hom. de degré 1 :
\[
(6.2.8)\qquad \beta : \mathbb{R}\Gamma(V_\eta, \Lambda_\nu) \longrightarrow
\mathbb{R}\Gamma(V_s, \Lambda_\nu)(-1)[-1] .
\]

\struck{\ill{}} Cet homomorphisme peut s'interpréter, par la proposition suivante, qui en
\uncertain{même} temps donne un \struck{\ill{}} calcul de $\mathbb{R}\Gamma(V, \Lambda_\nu)$
en termes de la donnée fondamentale \struck{\ill{}} (6.2.1) :

\emph{Proposition} 6.3. a) On a
\struck{$\underline{H}^i_{V_s}(\Lambda_{\nu\,V}) = 0$ pour $i \neq 2$,}
\struck{et $\underline{H}^2_{V_s}(\Lambda_{\nu\,V}) \simeq \Lambda_{\nu\,V_s}(-1)$,}
\[
(6.3.1)\qquad \underline{H}^i_{V_s}(\Lambda_{\nu\,V}) \simeq
\begin{cases} 0 & \text{si } i \neq 2 \\ \Lambda_{\nu\,V_s}(-1) & \text{si } i = 2 , \end{cases}
\]
en d'autres termes
\[
(6.3.2)\qquad \mathbb{R}\underline{\Gamma}_{V_s}(\Lambda_{\nu\,V}) \simeq \Lambda_{\nu\,V_s}(-1)[-2] ,
\]
d'où un isom.
\[
(6.3.3)\qquad \mathbb{R}\Gamma_{V_s}(V, \Lambda_\nu) \simeq \mathbb{R}\Gamma(V_s, \Lambda_\nu)(-1)[-2] .
\]
\note{« canonique » est ajouté sous la ligne, avec un renvoi à la fin de (6.3.3).}

b) Considérons le triangle

\begin{tikzcd}[column sep=small]
 & \mathbb{R}\Gamma(V_\eta, \Lambda_\nu) \arrow[dl, "\partial"'] & \\
\mathbb{R}\Gamma(V_s, \Lambda_\nu)(-1)[-2] \simeq \mathbb{R}\Gamma_{V_s}(\Lambda_{\nu\,V}) \arrow[rr, "i_*"] & & \mathbb{R}\Gamma(V, \Lambda_\nu) \arrow[ul, "j^*"']
\end{tikzcd}

(6.3.4).\note{L'isomorphisme du terme de gauche est glosé « par (6.3.3) ». La flèche horizontale
est notée d'un symbole lu $i_*$, peut-être $\iota_*$.}

Alors la flèche $\partial$ dudit triangle s'identifie
[via l'isomorphisme (6.3.3)] à $\beta$. \struck{\ill{}}
\add{À des signes et décalages près,}
\marginal{signe ?}

\emph{\struck{$\mathbb{R}\Gamma(V, \Lambda_\nu)$} Corollaire} 6.4
\struck{$\mathbb{R}\Gamma(V, \Lambda_\nu)$ s'identifie au isomorphe} Le triangle (6.3.4)
n'est autre que le triangle distingué \struck{construit avec} du « \uncertain{cylindre}
\uncertain{d'application} » du $\beta$ de (6.2.8).
\note{Le mot entre guillemets est surchargé ; « cône » ou « cylindre d'application » sont
possibles. La parenthèse « À des signes et décalages près » est ajoutée au-dessus de la ligne
biffée et porte sur 6.4.}

\page{159}
\note{Pagination de l'auteur « V.32 ».}

\emph{Démonstration} \uncertain{lem.}

\struck{$\mathbb{R}g_*(\Lambda_{V_{\bar\eta}})$}
\struck{$\simeq \mathbb{R}j_*\,\mathbb{R}h_*(\Lambda_{V_{\bar\eta}})$}
\struck{et \uncertain{introduisant} \ill{} $V_s \to V$, et \ill{} \ill{}}
\struck{$\mathbb{R}g_*(\Lambda_{V_{\bar\eta}})$}
\struck{$\simeq \mathbb{R}j_*\,\mathbb{R}h_*(\ldots)$}
\note{Le début de la démonstration, avec un premier diagramme triangulaire
($V_{\bar\eta} \to V$ notée $g$, $V_{\bar\eta} \to V_\eta$ notée $h$, $V_\eta \to V$ notée $j$),
est encadré et biffé ; une glose « $h_*(\Lambda_{V_{\bar\eta}})$ » coiffe le terme
$\mathbb{R}h_*(\Lambda_{V_{\bar\eta}})$.}

\begin{tikzcd}
V_{\bar s} \arrow[r, "\bar i"] \arrow[d, "g_0"'] & \bar V \arrow[d, "g"] & V_{\bar\eta} \arrow[l, hook', "\bar j"'] \arrow[d, "h"] \\
V_s \arrow[r, "i"'] & V & V_\eta \arrow[l, "j"]
\end{tikzcd}

\[
\begin{aligned}
\mathbb{R}j_*(\Lambda_{V_\eta}) &\simeq \mathbb{R}j_*\,\mathbb{R}\Gamma^\pi\,\mathbb{R}h_*(\Lambda_{V_{\bar\eta}}) \\
&\simeq \mathbb{R}\Gamma^\pi\,\mathbb{R}j_*\,\mathbb{R}h_*(\Lambda_{V_{\bar\eta}}) \\
&\simeq \mathbb{R}\Gamma^\pi\,\mathbb{R}g_*\,\mathbb{R}\bar j_*(\Lambda_{V_{\bar\eta}}) \\
&\xrightarrow{\ \sim\ } \mathbb{R}\Gamma^\pi\,\mathbb{R}g_*(\Lambda_{\bar V})
\end{aligned}
\]
car $\Lambda_{\bar V} \simeq \mathbb{R}\bar j_*(\Lambda_{V_{\bar\eta}})$, d'où
\[
\begin{aligned}
i^*\mathbb{R}j_*(\Lambda_{V_\eta}) &\simeq \mathbb{R}\underline{\Gamma}^\pi\, i^*\mathbb{R}g_*(\Lambda_{\bar V}) \\
&\simeq \mathbb{R}\underline{\Gamma}^\pi\,\mathbb{R}g_{0*}(\Lambda_{V_{\bar s}}) \\
&\simeq \mathbb{R}\underline{\Gamma}^\pi(\Lambda_{V_s})^{\mathrm{triv}} \\
&\simeq \Lambda_{V_s} \oplus \Lambda_{V_s}(-1)[-1]
\end{aligned}
\]
\note{Sur la page, le membre de gauche est écrit « $i^*\mathbb{R}j_*(V_{V_\eta})$ », sans doute
pour $i^*\mathbb{R}j_*(\Lambda_{V_\eta})$ ; on a normalisé, et de même un $\Lambda$ récrit
sur une autre lettre à la deuxième ligne.}
d'où aussitôt 6.3 a), en utilisant les relations habituelles entre les
$\underline{H}^i_{V_s}(\Lambda_V)$ et les $\mathbb{R}^i j_*(\Lambda_{V_\eta})$.

Compatibilité laissée au rédacteur en forme !

\page{160}
\note{Pagination de l'auteur « V.33 ». Le paragraphe 6.5 s'interrompt en bas de page et continue au-delà de ce lot.}

6.5. Utilisons maintenant les résultats du §~2, qui indiquent,
\struck{si $n = \dim \mathcal{O}_{X_0,x}$} posant
\[
(6.5.1)\qquad n = \dim \mathcal{O}_{X_0,x}
\]
que l'on a
\[
(6.5.2)\qquad H^i(V_{\bar\eta}) = 0 \quad \text{si } i > n .
\]
\add{Utilisons les} \struck{Compte tenu} des \struck{la} suites exactes déduites
\struck{de 6.3.4} de la suite spectrale de Hochschild--Serre
\struck{$H^i(V) \to H^i(V_\eta)$}
\[
H^*(V_\eta) \Longleftarrow E_2^{pq} = H^p(\pi, H^q(V_{\bar\eta}))
\]
[déduite de $\mathbb{R}\Gamma(V_\eta) \simeq \mathbb{R}\Gamma^\pi\,\mathbb{R}\Gamma(V_{\bar\eta})$] :
\[
(6.5.3)\qquad \cdots \longrightarrow H^{i-1}(V_{\bar\eta})_\pi \longrightarrow H^i(V_\eta)
\longrightarrow H^i(V_{\bar\eta})^\pi \longrightarrow 0 ,
\]
on trouve donc
\[
(6.5.4)\qquad H^i(V_\eta) = 0 \quad \text{si } i > n+1 .
\]
\struck{Utilisons} Écrivons la suite exacte déduite de 6.3.4
\[
(6.5.5)\qquad \longrightarrow H^{i-2}(V_s)(-1) \longrightarrow H^i(V) \longrightarrow H^i(V_\eta)
\longrightarrow H^{i-1}(V_s) \longrightarrow H^{i-1}(V) \longrightarrow \cdots ,
\]
\note{Dans (6.5.5), entre $H^{i-2}(V_s)(-1)$ et $H^i(V)$, un terme est noirci et illisible ;
dans (6.5.3), l'exposant du premier terme est écrit sur un autre.}
on \struck{trouve} \uncertain{voit} que 6.5.4 s'exprime aussi par les relations :
\[
(6.5.6)\qquad i_* : H^i(V, \Lambda) \longrightarrow H^{i-2}(V_s)
\]
est un isom. si $i > n$\ill{}, un épimorphisme si $i = n$\ill{}.
\note{Les bornes de (6.5.6) sont surchargées : on lit « $i > n$ » suivi d'un chiffre noirci,
et « $i = n$ » suivi d'un autre, que l'on ne restitue pas. La flèche $i_*$ va, sur la page,
de $H^i(V, \Lambda)$ vers $H^{i-2}(V_s)$, sens contraire de celui de (6.5.5).}

\drawing{Seul en bas de page, un « 6 » isolé, sans suite.}

\end{document}
